%! TEX root = ketentuan_riset.tex
% the main TeX file which is intended to compile, :VimtexReload after adjustment
   % this file should be included in the main file
% :h vimtex-tex-root
\documentclass[../riset_senja.tex]{subfiles}
%ensure the class of docs

\begin{document}

\section{Persyaratan}%
\label{sec:Persyaratan}

\noindent\textbf{Petunjuk Umum:} 1500 -- 2500 kata.

\subsection{A. Judul Penelitian}
Tuliskan judul penelitian.

\subsection{B. Latar Belakang}
Uraikan secara singkat topik isu yang ingin Anda teliti dan mengapa signifikan untuk Anda teliti.

\subsection{C. Perumusan Permasalahan (\textit{Statement of Problem})}
Uraikan secara singkat apa yang Anda ketahui tentang topik isu tersebut dan jelaskan secara ringkas mengapa masih perlunya Anda teliti. Tunjukkan bahwa solusi yang ada terhadap isu tersebut masih belum sepenuhnya menyelesaikan permasalahan, sehingga Anda ingin melakukan penelitian.

\subsection{D. Pertanyaan/Tujuan Penelitian}
Rumuskan tujuan pertanyaan penelitian.

\subsection{E. Kelogisan (\textit{Rationale})}
Jelaskan bagaimana pertanyaan penelitian mendukung topik isu besar yang diangkat dalam latar belakang penelitian. Jelaskan hipotesis (jika ada) dan/atau model penelitian yang mendukung tujuan/pertanyaan penelitian. Jelaskan pula kontribusi teoritis dan praktis jika hipotesis tidak terbukti.

\subsection{F. Metode dan Desain}
\begin{enumerate}[label=\alph*.]
	\item Jelaskan bagaimana Anda akan mengumpulkan data dan mengapa.
	\item Jelaskan mengapa metode ini adalah terbaik untuk mencapai tujuan Anda.
	\item Jelaskan analisis dan hasil yang mendukung maupun tidak mendukung hipotesis.
	\item Cantumkan outline jadwal penelitian dari awal sampai selesai.
\end{enumerate}

\subsection{G. Signifikansi/Manfaat}
Uraikan secara umum, bagaimana penelitian yang Anda usulkan dapat berguna baik secara teoritis maupun praktis.

\subsection{H. Daftar Pustaka}


\section{Pendahuluan}%
\label{sec:Pendahuluan}


Uraikan secara umum, bagaimana penelitian yang Anda usulkan dapat berguna baik secara teoritis maupun praktis.






\bibliographystyle{apalike}
\bibliography{../filename.bib} % required for citecompletion, not work in mainfile
\end{document}

